\section{Ablation Studies}

\subsection{Prompt-Type Training}
\label{sec:ablation_prompt}

To validate that prompt robustness arises from mixed-prompt training rather than being an inherent property of fine-tuning, we train three decoder variants: point-only, box-only, and mixed (default).

\begin{table}[t]
\centering
\caption{\textbf{Prompt-type ablation} (500 test images). Mixed training produces the most robust decoder. Box-only training fails catastrophically on point prompts.}
\label{tab:ablation_prompt}
\small
\begin{tabular}{lcccc}
\toprule
\multirow{2}{*}{Training} & \multicolumn{4}{c}{Test Prompt mIoU} \\
\cmidrule(lr){2-5}
 & CoM & Edge & Grid & Random \\
\midrule
Point-only & .699 & .672 & .725 & .768 \\
Box-only & .376 & .501 & .389 & .455 \\
\textbf{Mixed (ours)} & \textbf{.738} & \textbf{.694} & \textbf{.756} & \textbf{.774} \\
\bottomrule
\end{tabular}
\end{table}

Mixed-prompt training outperforms point-only on all test strategies (by 2--4 mIoU on CoM, Edge, and Grid; 0.6 on Random where point-only already performs well) and box-only by 19--37 mIoU. Box-only training scores 0.376 mIoU on CoM (near-catastrophic), showing that prompt robustness requires explicit exposure to both prompt formats during training, not just fine-tuning per se.

\subsection{Loss Function}

\begin{table}[t]
\centering
\caption{\textbf{Loss function ablation} (500 test images, Center-of-Mass). All variants perform similarly; Dice-heavy is marginally best.}
\label{tab:ablation_loss}
\small
\begin{tabular}{lccccc}
\toprule
Loss Config & mIoU & Dice & $S_\alpha$ & $F_\beta^w$ & MAE$\downarrow$ \\
\midrule
BCE only (1.0/0.0) & .713 & .796 & .823 & .782 & .042 \\
Dice only (0.0/1.0) & .712 & .793 & .815 & .778 & .046 \\
BCE-heavy (0.7/0.3) & .708 & .789 & .818 & .781 & .044 \\
Dice-heavy (0.3/0.7) & \textbf{.729} & \textbf{.810} & \textbf{.829} & \textbf{.801} & \textbf{.043} \\
\midrule
Ours (0.5/0.5)$^\dagger$ & .738 & .819 & .847 & .819 & .028 \\
\bottomrule
\multicolumn{6}{l}{\footnotesize $^\dagger$Full 2,026-image evaluation; ablation rows use 500 images.}
\end{tabular}
\end{table}

All four loss variants fall within $\pm$2.1 mIoU of each other, so the decoder is not sensitive to loss weighting. The Dice-heavy variant is marginally best, hinting that region-level supervision is slightly more useful for COD than pixel-level BCE.

\subsection{LLPM Component Analysis}

To validate that both LLPM branches contribute, we train variants with individual branches disabled.

\begin{table}[t]
\centering
\caption{\textbf{LLPM component ablation} (500 test images, Edge prompt mIoU). Both branches and the learnable gate contribute; removing any component degrades performance.}
\label{tab:ablation_llpm}
\small
\begin{tabular}{lccc}
\toprule
LLPM Variant & Edge mIoU & $\Delta$ & $\alpha_{\text{final}}$ \\
\midrule
No LLPM (decoder-only) & .694 & --- & --- \\
Edge branch only & .713 & +1.9 & 0.091 \\
Enhancement only & .713 & +1.9 & 0.211 \\
No gate ($\alpha\!=\!1$) & .709 & +1.5 & 1.0 (fixed) \\
\textbf{Full LLPM (ours)} & \textbf{.719} & \textbf{+2.5} & 0.185 \\
\bottomrule
\end{tabular}
\end{table}

Each branch alone gains +1.9 Edge mIoU, but the full LLPM reaches +2.5, so the two branches are not redundant. The learnable gate matters too: fixing $\alpha\!=\!1$ gives the weakest variant (+1.5), since the model can no longer regulate enhancement strength. The edge-only variant's $\alpha$ barely moves from its initialization (0.091 vs.\ 0.185 for the full LLPM), suggesting the edge attention map has limited standalone value without the enhancement branch to act on it.

\subsection{Data Augmentation}

We test progressively richer augmentation beyond horizontal flips: adding random rotation ($\pm$5\textdegree) and pixel shift ($\pm$16px, half a ViT patch).

\begin{table}[t]
\centering
\caption{\textbf{Augmentation ablation} (500 test images). Richer augmentation slightly hurts CoM/Grid/Random but improves Edge, the hardest prompt. Differences are small ($<$2 mIoU), indicating robustness to augmentation choice.}
\label{tab:ablation_aug}
\small
\begin{tabular}{lcccc}
\toprule
\multirow{2}{*}{Augmentation} & \multicolumn{4}{c}{Test Prompt mIoU} \\
\cmidrule(lr){2-5}
 & CoM & Edge & Grid & Random \\
\midrule
\textbf{Flip only (ours)} & \textbf{.738} & .694 & \textbf{.756} & \textbf{.775} \\
Flip + rotation & .722 & .708 & .741 & .768 \\
Flip + rot.\ + shift & .724 & \textbf{.718} & .738 & .765 \\
\bottomrule
\end{tabular}
\end{table}

Adding rotation and pixel shift increases training samples (9,120 and 12,160 vs.\ 6,080) and lowers training loss (0.068 and 0.064 vs.\ 0.072), but does not improve test mIoU on most prompts. The exception is edge prompts, where richer augmentation gains +2.4 mIoU (flip+rot+shift vs.\ flip-only), likely because geometric perturbations expose the decoder to more boundary configurations. Flip-only is near-optimal overall.

\subsection{Qualitative Results}

\begin{figure*}[t]
    \centering
    \includegraphics[width=\textwidth]{figures/qualitative_comparison.png}
    \caption{\textbf{Qualitative comparison.} Each row: input, GT mask, base SAM, our decoder (center-of-mass prompts). Top: hard cases where base SAM fails catastrophically. Bottom: easier cases where both succeed. The specialized decoder consistently follows true boundaries even when coloration matches the background.}
    \label{fig:qualitative}
\end{figure*}

Base SAM tends to segment only the most salient patch near the prompt, missing large portions of camouflaged objects. The specialized decoder follows object boundaries even where the animal's coloration is nearly identical to the background (\cref{fig:qualitative}).
